%% kapa: exam=1 round=20 year=2009 period=1 form=A subject=civil_law subject_ko=민법 seq=26 number=26
%% source: Q-Net 감정평가사 기출문제 (cst00309) artlSeq=5024178 file=extracted/1008/5024178/2009년 제20회 감정평가사 1차시험 기출문제/1교시.hwp
%% conversion: kordoc 4.24.51 -> tools/gen1.py (auto); review=proofread-sonnet; flags=none
\begin{question}{26}
다음 (가), (나)에 관한 설명으로 옳지 \uline{않은} 것은? (다툼이 있으면 판례에 의함)
\begin{framedtext}
(가) 甲종중은 건물을 신축하여 종원(宗員) 乙과의 명의신탁 약정에 따라 乙명의로 소유권보존등기를 하였다(위법한 목적이 없음을 전제로 함). \par
(나) 乙이 위 건물을 丙에게 매도하고 소유권이전등기까지 경료해 주었다. \par
\end{framedtext}
\begin{choices}
  \item (가)의 경우, 甲은 명의신탁약정을 해지하고 乙에게 건물에 대한 소유권이전등기를 청구할 수 있다.
  \item (가)의 경우, 건물의 소유권은 대외적으로 乙에게 귀속된다.
  \item (나)의 경우, 丙이 甲ㆍ乙 사이의 명의신탁에 관하여 악의라면 건물의 소유권을 취득하지 못한다.
  \item (나)의 경우, 甲은 丙에게 그 건물의 반환과 이전등기의 말소를 직접 청구할 수 없다.
  \item (나)의 경우, 丙이 乙의 배임행위에 적극 가담한 경우에는 乙과 丙의 계약은 반사회적 법률행위로서 무효이다.
\end{choices}
\end{question}
\answer{3}
